
2.8.
Soit @@M0@@ un complexe parfait sur @@M1@@. Il existe un entier @@M2@@ tel que @@M3@@ soit un complexe @@M4@@-parfait, et que @@M5@@ soit isomorphe à @@M6@@. Si on applique 2.7 au complexe @@M7@@, la suite des triangles distingués construite donne en appliquant le foncteur @@M8@@ une suite de triangles distingués de @@M9@@ (il est clair que @@M10@@ transforme complexe @@M11@@-parfait en complexe parfait). Par ailleurs, on a :
@@M12@@
(notons que @@M13@@ d'après la platitude de @@M14@@) pour @@M15@@ dans @@M16@@.

On obtient donc dans @@M17@@ l'égalité :
@@M18@@
où les @@M19@@ sont les complexes parfaits sur @@M20@@ définis dans 2.7. Il en résulte que @@M21@@ est engendré comme @@M22@@-module par les puissances de @@M23@@.

Pour voir que @@M24@@ sont linéairement dépendants, on considère l'homomorphisme~: @@M25@@ qui à la classe d'un faisceau localement libre de type fini @@M26@@ associe la classe du complexe égal à @@M27@@ en degré @@M28@@, et nul pour les autres degrés. L'équation (1.9.1) donne alors une relation de dépendance dans @@M29@@.

Pour voir que @@M30@@ sont linéairement indépendants, on peut utiliser le morphisme @@M31@@-linéaire @@M32@@ (IV 2.11.1). On a :
@@M33@@
@@M34@@
Si on a une relation du type :
@@M35@@
on la multiplie par @@M36@@, et on applique @@M37@@; on trouve alors @@M38@@; on 
